\documentclass[11pt]{article}
\usepackage[margin=1in]{geometry}
\usepackage{amsmath, amssymb}

\title{How Hard Is It to Factor Homogeneous Polynomials?}
\author{ScottLean4 \textperiodcentered\ for Dana Scott}
\date{}

\begin{document}
\maketitle

The answer is a pleasant surprise: \emph{computationally, it is easy---polynomial
time.} This is worth savoring, because the sibling problem, \emph{integer}
factorization, is believed hard (no polynomial algorithm is known, and much of
public-key cryptography rests on that). Polynomials factor in polynomial time;
integers, apparently, do not.

\section*{Binary forms (two variables): as easy as univariate}

A binary form $F(x,y)$ of degree $d$ is equivalent to a one-variable polynomial.
\textbf{Dehomogenize}---set $y = 1$ to get $f(x) = F(x,1)$ of degree $\le d$---%
factor that, then \textbf{re-homogenize} each factor (restoring any lost factor
of $y$). So factoring binary forms is \emph{equivalent} to univariate
factorization:
\begin{itemize}
  \item Over $\mathbb{Q}$: polynomial time, by \textbf{Lenstra--Lenstra--Lov\'asz
    (LLL, 1982)}---lattice basis reduction made univariate rational factoring
    provably efficient.
  \item Over a finite field $\mathbb{F}_q$: polynomial time (randomized), via
    \textbf{Berlekamp} and \textbf{Cantor--Zassenhaus}.
\end{itemize}

Geometrically, a binary form of degree $d$ is exactly $d$ points on the
projective line $\mathbb{P}^1$ (its roots with multiplicity). Over an
algebraically closed field it splits completely into $d$ linear factors. The
degenerate conic
\[
  x^2 - y^2 = (x - y)(x + y)
\]
is just ``the binary quadratic has two roots on $\mathbb{P}^1$''---the pair of
lines $x = y$ and $x = -y$.

\section*{Ternary and higher forms: still polynomial time}

A homogeneous polynomial in $n$ variables dehomogenizes to a general polynomial
in $n-1$ variables. Multivariate factorization over $\mathbb{Q}$ or
$\mathbb{F}_q$ is \emph{also} polynomial time---\textbf{Kaltofen (1980s)}---by
evaluating to the univariate case and then \textbf{Hensel lifting} to recover the
multivariate factors. Homogeneity is a mild convenience (each factor is itself
homogeneous), not a change in complexity class.

\section*{Where the difficulty actually lives}

The hardness is in the constants and the field, not the complexity class.
\begin{enumerate}
  \item \textbf{The base field decides the answer's shape.} Over $\mathbb{C}$
    everything splits into linear forms; over $\mathbb{R}$ one gets linear and
    irreducible quadratic forms; over $\mathbb{Q}$ higher-degree irreducibles
    persist and factor coefficients may live in extensions. Same algorithm, very
    different output.
  \item \textbf{Coefficient growth.} LLL and Hensel lifting are polynomial, but
    intermediate coefficient heights can swell, so a large-height form is slow in
    practice though poly-time in theory.
  \item \textbf{Verify versus discover.} \emph{Verifying} a factorization is
    trivial (multiply out). \emph{Discovering} it is the algorithmic work above.
    Both are easy in the complexity sense; only the second is interesting.
\end{enumerate}

\section*{The formalization angle}

So the \emph{algorithm} is easy---but a \emph{certified} factoring procedure is a
different matter. Mathlib has the \emph{theory} (polynomials over a field form a
unique factorization domain, irreducibility, and so on), but there is no standard
\emph{tactic} that returns the factors \emph{with a kernel-checked proof} the way
\texttt{ring} verifies a guess. That gap is engineering, not complexity:
producing the factors is fast; producing them together with a formal certificate
is the work no one has fully packaged.

\end{document}
